% Translated from sections/06-discussion.tex on 2026-09-23. The English file is
% canonical. Change it first, then carry the change here.
\section{Discussion et limites}
\label{sec:errors}
% Both labels point here, as in the English file.
\label{sec:limitations}

\subsection{La règle d'acceptation, face à la littérature et à ses propres failles}

La littérature sur la correction a posteriori converge vers un même
avertissement. Un correcteur LLM dégrade une reconnaissance déjà
bonne~\cite{kanerva2025nofreelunches}, et la surcorrection des entrées peu
bruitées est le problème récurrent de HIPE-OCRepair~\cite{ehrmann2026hipeocrepair}.
Les deux décrivent un correcteur qui réécrit quand il devrait s'abstenir.
L'ancrage par marqueurs transforme ce problème de qualité en problème
détectable. Les \corrRewritesCaught{}~refus sont des réponses qui ont modifié
le texte hors des marqueurs, et qu'une règle sans ancrage aurait acceptées sans
bruit. Mais cette branche se déclenche sur la moindre différence d'octet. Elle
majore donc la réécriture fluide sans la compter. Une espace déplacée et une
proposition réécrite tombent dans le même panier, et seuls
\corrRewritesSubstantial{} des \corrRewritesCaught{}~refus déplaçaient plus de
deux caractères.

\textbf{Une faille dans la règle.} Elle n'est pas aussi étanche que la
section~\ref{sec:system} le décrit. \texttt{suspect\_spans} situe chaque mot
suspect par son décalage en caractères quand le texte brut du reconnaisseur et
celui de la page s'alignent. Sinon, elle cherche le mot et retient sa
\emph{première} occurrence. Deux occurrences de faible confiance d'un même mot
se ramènent alors au même indice. La ligne reçoit deux paires de marqueurs
autour d'un seul mot, le modèle voit une ligne corrompue, et la règle compare
la réponse à cette même construction corrompue, puis l'accepte. La règle ne
peut pas détecter ce cas, parce que la corruption a lieu avant son exécution.
Et cette branche n'a rien de marginal, puisque \corrOnFindFallback{} des
\corrOffered{}~lignes soumises l'ont empruntée. C'est la limite la plus grave
de ce rapport. Le taux d'acceptation mesuré majore le comportement sain de la
règle, et \textbf{C2 n'est étayée que sur moins d'un tiers des lignes soumises
par E2}. Une réserve sur ce chiffre s'impose. L'échec du contrôle par décalage
est nécessaire au défaut mais pas suffisant, car il faut aussi deux
occurrences de faible confiance d'un même mot sur une même ligne.
Le taux de \corrFindFallbackShare{} borne donc l'exposition au lieu d'estimer le dommage.
La correction relève de la chaîne de traitement. Il faut chercher à partir de
la fin du segment précédent, et refuser de marquer une ligne dont les segments
se chevauchent.

\textbf{Et E2 mesure le refus, non l'amélioration.} Savoir si les
modifications acceptées ont amélioré le texte exige une référence pour les
pages tirées, et le jeu de référence n'en couvre aucune. Son échantillon
compte en outre \corrPages{} des \corrReachPagesWithSpan{}~pages qui portent un
segment suspect. C'est un échantillon du corpus, et non le corpus.

\subsection{La localisation, et la structure}

L'exactitude est de \locLivreAccuracyWorkbook{} au niveau du Livre et de
\locChapterAccuracyWorkbook{} au niveau du chapitre, sur les mêmes pièces. Le
préenregistrement avait prévu ce cas en ces termes: «\textenglish{\emph{Correct
CHAPTER far below correct LIVRE points at the span, not the search}}».\footnote{Un
taux de chapitre exact très inférieur au taux de Livre exact met en cause
l'étendue retenue, et non la recherche.} Il prescrivait du même souffle de
régler \texttt{CHAIN\_GAP\_PENALTY} et \texttt{WINDOW\_SLACK} avant de
conclure. Ce
réglage n'a pas été fait. La lecture par l'étendue est donc l'attente
préenregistrée confirmée dans son sens, et non une cause éprouvée. Les
dénominateurs disent ce que les taux ne disent pas. Seules
\locChapterScoreableWorkbook{} des \locPiecesWorkbook{}~pièces certaines portent
un numéro de chapitre imprimé déchiffrable. L'exactitude au niveau du chapitre
est donc le plus faible des deux chiffres, par les données comme par la
valeur.

Sur la structure de la famille~A, \structSpuriousA{}~débuts de chapitre
proposés à tort subsistent contre \structMissingA{}~manques, après qu'un
lecteur a fait 128~promotions et 217~rétrogradations à la main. Ce résidu est
une surproposition, et il ne renseigne pas sur le détecteur seul. La passe
manuelle a supprimé une part inconnue des manques, et en borner l'effet place
le rappel du détecteur seul à 77,0\,\% au plus bas, contre les
\structRecallHandCheckedA{} rapportés.

\subsection{Ce qui manque}

\textbf{Trois des cinq expériences ne sont pas mesurées.} E1, E3 et E4 ne sont
plus bloquées, puisque le jeu de référence existe avec \goldPages{}~pages. Mais
elles n'ont pas été exécutées avant la clôture de ce rapport. La conséquence
reste la même, et elle est lourde. \textbf{L'affirmation~C1 n'est pas évaluée.}
Ce rapport montre que la chaîne fonctionne, à quelle vitesse, ce que sa règle
d'acceptation refuse et avec quelle justesse elle localise un passage. Il ne
montre pas que le modèle affiné lit mieux que les modèles qu'il a remplacés.
Le projet affirme depuis un an une table de confusion où le \qtr{ſ} long se
confond avec \qtr{f}, \qtr{u} avec \qtr{n} et \qtr{c} avec \qtr{e}. Elle est
désormais calculable, et toujours pas calculée.

\textbf{Le jeu de référence est postédité et n'a qu'un annotateur.} Ses pages
ont été corrigées sur la sortie d'un modèle au lieu d'être transcrites à
l'aveugle. Il penche donc vers le modèle affiné, dans le sens qui
flatterait C1. Et aucune page n'a été faite deux fois, si bien qu'il n'existe
pas de plancher en dessous duquel une différence serait écartée. La
section~\ref{sec:protocol} énonce ces deux points. Un second annotateur et une
passe à l'aveugle sont les deux premières choses qu'un successeur devrait
ajouter, dans cet ordre. En dessous de 30~cas dans une case, E3 et E4 se
réduisent à une ventilation descriptive, et certaines cases tomberont sous ce
seuil.

\textbf{La validation n'est pas un test.} Les \trainPagesVal{}~pages réservées
ont servi à choisir le point de contrôle livré. Elles ont été séparées page par
page à partir d'un seul volume, si bien que les mêmes formes et la même usure
des caractères apparaissent des deux côtés. Le modèle a aussi lu une partie du
texte sur lequel il est évalué. Aucune page d'\emph{Amadis de Gaule} ne figure
dans la liste d'entraînement, mais le \emph{Trésor} reprend des extraits
d'\emph{Amadis}. L'ampleur de cet effet ne peut être estimée sans savoir quels
passages T.1 reprend.

\textbf{Transkribus est un système de comparaison confondu.} Il est rapporté
comme tel au lieu d'être écarté, car l'écarter masquerait la comparaison qu'un
lecteur souhaite le plus. Sa propre sortie a été cherchée sans être trouvée,
si bien que le quatrième système de E1 pourrait manquer entièrement.
\textbf{Cinq des six seuils ne sont pas calibrés} (tableau~\ref{tab:thresholds}),
pas plus que les deux que le préenregistrement pose comme conditions à la
lecture du résultat par chapitre. \textbf{Enfin, les durées ne couvrent que le
traitement par lots.} Le chiffre de \ocrSecondsPerPage{}~s par page vaut
pour ce programme sur ce GPU. Ce n'est pas une affirmation sur le système déployé,
dont la base d'exécution n'a pas été lue.
